\documentclass[12pt]{article}
\usepackage[margin=1in]{geometry}
\usepackage{kotex}

\setlength{\parindent}{0pt}
\setlength{\parskip}{1\baselineskip}

\title{DISTANCE v0.1 \\[0.5em]
\large A Philosophical Framework for the Direction of the Universe \\
}
\author{Kangbin Yim\\Soonchunhyang University}
\date{2024}

\begin{document}
\maketitle

우리는 오랫동안 우주를 공간 속에 놓인 물체들의 집합으로 이해해 왔다. 별은 공간 안에 떠 있고, 행성은 그 사이를 움직이며, 생명은 특정한 장소에서 태어나고 사라진다고 생각해 왔다. 우리는 시간을 흐르는 것으로 느끼고, 공간을 비어 있는 배경처럼 인식하며, 거리와 속도를 통해 세계를 설명해 왔다. 이러한 감각은 너무 자연스럽기 때문에 대부분의 사람들은 그것을 의심하지 않는다.

뉴턴은 시간과 공간을 절대적인 스케일로 보았고, 아인슈타인은 그것을 상대적 구조로 바꾸었다. 현대물리학은 우주를 설명하기 위해 수많은 수식과 모델을 만들어 냈고, 인간은 그 체계를 이용해 문명을 발전시켜 왔다. 그러나 그 모든 변화 속에서도 우리는 여전히 시간과 공간 자체를 우주의 가장 근본적인 실체처럼 받아들여 왔다. 우리는 사건이 시간 속에서 발생하고, 물체가 공간 속에서 움직인다고 믿어 왔다.

하지만 어느 순간부터 나는 조금 다른 질문을 하기 시작했다.

정말 우주의 가장 근본적인 실체는 시간과 공간일까.
정말 사건은 시간 속에서 발생하고, 물체는 공간 속에서 움직이는 것일까.
아니면 시간과 공간은 어떤 더 근본적인 작용을 인간이 이해하기 쉽게 번역한 인터페이스에 불과한 것은 아닐까.

우리는 시계의 초침이 움직이는 것을 보며 시간이 흘렀다고 말한다. 그러나 실제로 일어난 것은 에너지 전달이다. 태엽이 풀리고, 배터리의 화학 반응이 전류를 만들고, 회로가 작동하며, 그 결과 바늘의 위치가 변한다. 우리가 시간이라고 부르는 것은 그 변화 이후에 인간이 붙인 해석일 뿐일 수 있다.

공간 역시 마찬가지다. 우리는 어떤 물체가 이곳에서 저곳으로 이동했다고 말한다. 그러나 실제로는 에너지 상태가 변화했고, 어떤 관계가 재배치되었으며, 인간은 그것을 거리와 위치라는 형식으로 번역했을 뿐인지도 모른다.

나는 점점 우주를 "에너지 상태들의 관계 변화"로 보기 시작했다.

어떤 존재가 다른 존재와 서로 다른 에너지 상태를 가지게 되면, 그 사이에는 변화 가능성이 생긴다. 에너지는 완전히 균일한 상태에 머무르지 않고, 더 낮은 상태를 향해 흐르려는 경향을 가진다. 나는 이러한 흐름의 가능성과 방향성, 그리고 관계 변화의 잠재성을 "모멘텀"이라고 부르고자 한다.

여기서 말하는 모멘텀은 단순한 물리학적 운동량 개념이 아니다. 그것은 존재 간 에너지 상태 차이로 인해 발생하는 관계 변화 가능성에 가깝다. 우주의 모든 변화는 이 모멘텀으로부터 시작된다.

그리고 나는 바로 이 모멘텀 관계 상태를 "디스턴스"라고 생각하게 되었다.

여기서 말하는 디스턴스는 단순한 공간적 거리나 시각적 단절 상태가 아니다. 디스턴스는 에너지 상태 차이로 인해 형성되는 비대칭적 관계 상태이다. 서로 다른 에너지 수준을 가진 존재들 사이에서 에너지 전달의 방향성과 강도가 형성되며, 이들이 서로 관계하게 되는 상태 자체이다.

모멘텀이 크면 디스턴스는 작고, 모멘텀이 작으면 디스턴스는 크다.

이 관점에서 우주는 공간 속 물체들의 집합이 아니라 디스턴스들의 장이 된다. 에너지 상태 차이들이 모멘텀을 만들고, 모멘텀들이 관계를 만들며, 그 관계들이 끊임없이 재배치되는 거대한 흐름이다.

시간은 그 흐름을 인간이 인식하는 방식이며, 공간은 그 관계를 번역하기 위해 인간이 사용하는 인터페이스에 가깝다.

우리는 존재를 독립된 객체처럼 인식한다. 인간, 생명체, 별, 행성, 국가, 문명 모두 독립적으로 존재하는 실체처럼 느껴진다. 그러나 어쩌면 존재는 독립된 실체가 아니라 특정 수준의 에너지 상태가 잠시 만들어낸 상(image)에 가까운 것인지도 모른다.

존재는 고정된 물체가 아니라 디스턴스를 변화시키는 과정 속에서 잠시 유지되는 흐름 상태일 수 있다. 그리고 그 존재들은 다시 다른 존재들과 상호작용하며 서로의 디스턴스를 변화시킨다.

겉으로 안정적으로 보이는 존재조차 실제로는 다른 존재들의 디스턴스를 변화시키기 위해 끊임없이 작동하고 있는 복합체일 수 있다. 하나의 존재는 다른 존재들의 모멘텀을 변화시키고, 관계를 재배치하며, 더 높은 엔트로피 방향으로 우주를 움직이는 로컬 엔진처럼 작동한다. 그러나 동시에 그 존재 자신 역시 또 다른 존재들에 의해 끊임없이 변화당하고 있다.

우주는 고정된 존재들의 집합이 아니라 서로의 디스턴스를 변화시키는 과정들의 총합처럼 보인다.

생명 역시 예외가 아니다.

우리는 오랫동안 생명을 엔트로피에 저항하는 존재처럼 생각해 왔다. 생명체는 자기 내부의 질서를 유지하고 스스로를 복제하며 구조를 안정화하기 때문이다. 그러나 다른 관점에서 보면 생명은 우주적 흐름에 저항하는 존재가 아니라 디스턴스를 탐색하고 연결하며 재배치하는 자기조직화 구조이다.

식물은 태양과 행성 사이의 에너지 상태 차이를 이용하고, 동물은 화학 에너지 차이를 따라 움직이며, 미생물은 죽은 구조를 분해하여 새로운 흐름을 만든다. 생명은 우주 바깥의 기적이 아니라 디스턴스 흐름 위에서 형성된 국소적 소용돌이에 가깝다.

중요한 것은 생명 자체가 아니다. 중심은 언제나 흐름이다.

생명은 그 흐름이 특정 조건에서 취한 형식일 뿐이다. 은하의 회전도, 태풍의 소용돌이도, 별 내부의 플라즈마 흐름도, 생태계의 순환도 본질적으로는 모두 디스턴스 위에서 발생한 패턴들이다.

의식 역시 예외가 아닐 수 있다.

우리는 의식을 특별하고 초월적인 것으로 느낀다. 그러나 의식 또한 에너지 흐름 위에서 발생한 고도의 자기참조 구조일 수 있다. 신경망은 신호를 교환하고, 에너지를 소비하며, 내부 상태를 끊임없이 갱신한다. 생각과 기억과 감정은 고정된 물체가 아니라 흐름의 상태이다.

어쩌면 의식은 디스턴스 흐름이 자기 자신을 부분적으로 관찰하기 시작한 상태인지도 모른다.

그러나 그것이 인간을 우주의 중심으로 만드는 것은 아니다.

인간은 단지 디스턴스 흐름이 만들어낸 수많은 국소 구조 중 하나일 뿐이다. 인간 문명 역시 디스턴스 흐름이 특정 규모에서 만들어낸 하나의 복합 구조에 불과하다. 도시와 산업과 네트워크와 정보 시스템은 모두 막대한 에너지 상태 차이를 연결하고 재배치하는 구조들이다.

문명은 질서처럼 보이지만 동시에 더 큰 엔트로피와 디스턴스 변화를 만들어내는 흐름 증폭 장치이기도 하다.

이 관점에서 엔트로피 역시 다시 정의된다.

엔트로피는 단순한 무질서 증가가 아니다. 그것은 존재들 간 디스턴스가 점점 증가하는 방향성이다. 즉 모멘텀이 감소하고, 에너지 전달 가능성이 줄어들며, 관계 강도가 약해지는 방향이다.

우주의 마지막은 단순한 열적 죽음이 아니라 모든 존재 간 디스턴스가 무한대로 증가하여 더 이상 어떤 관계 변화도 발생하지 않는 상태일 수 있다. 최소 수준의 에너지 담지 주체들만 남고, 그 사이에 어떠한 에너지 이동도 존재하지 않는 진정한 정지 상태이다.

그리고 이 관점에서 빅뱅 역시 다시 보인다.

빅뱅은 단순히 공간이 팽창하는 사건이 아니라 극도로 낮은 디스턴스 상태에서 디스턴스가 증가하기 시작한 사건일 수 있다. 우주는 공간적으로 퍼져 나가는 것이 아니라 관계적으로 멀어지고 있는 것인지도 모른다.

그러나 여기에서 하나의 질문이 남는다.

만약 우주의 흐름이 궁극적으로 모든 디스턴스를 무한대로 증가시키는 방향이라면, 왜 우주는 다시 새로운 모멘텀을 만들어내는가. 왜 차이는 계속 생성되는가. 왜 우주는 완전한 정지 상태에 머무르지 않는가.

나는 아직 이 질문에 대해 완전한 답을 가지고 있지 않다.

어쩌면 우주의 극단적인 엔트로피 상태와 극단적인 저디스턴스 상태는 인간이 생각하는 것처럼 서로 완전히 분리된 양 끝이 아닐 수도 있다. 어쩌면 가장 멀어진 상태가 동시에 가장 가까운 상태와 이어져 있는 하나의 환형 구조일 수도 있다. 혹은 인간은 아직 우주를 설명하기 위한 올바른 질문 자체를 발견하지 못한 것인지도 모른다.

그러나 적어도 나는 우리가 너무 오랫동안 시간과 공간이라는 인터페이스를 우주의 실체 자체로 착각해 왔을 가능성은 충분히 존재한다고 생각한다.

DISTANCE는 새로운 절대적 우주론을 제시하려는 시도가 아니다.

오히려 이것은 인간이 너무 빨리 우주를 이해했다고 믿어 버린 방식에 대한 긴 의심이다.

우리는 시간과 공간을 너무 쉽게 절대적 실체로 받아들였고, 그 위에 수많은 과학과 철학과 사회 구조를 쌓아 올렸다. 하지만 어쩌면 우리는 처음부터 인터페이스를 실재로 착각하고 있었는지도 모른다.

그리고 이 이야기는 바로 그 의심으로부터 시작된다.


\end{document}